\section*{Chapter \ref{ch:b1:s-block} --- Hydrogen and the s-Block}

\begin{solution}{exo:b1:s-block:1}
\ce{KH} and \ce{CaH2}: saline; \ce{SiH4} and \ce{H2S}: covalent;
\ce{PdH_{0.6}}: metallic. \ce{CaH2 + 2H2O -> Ca(OH)2 + 2H2}.
\end{solution}

\begin{solution}{exo:b1:s-block:2}
\ce{2Li + 2H2O -> 2LiOH + H2}, and the same with Na and K;
\ce{Mg + 2HCl -> MgCl2 + H2}.
\end{solution}

\begin{solution}{exo:b1:s-block:3}
Basic: \ce{Na2O}, \ce{MgO} (\ce{MgO + 2HCl -> MgCl2 + H2O}). Acidic:
\ce{SiO2}, \ce{SO3}, \ce{CO2} (\ce{SO3 + H2O -> H2SO4}). Amphoteric:
\ce{Al2O3}, \ce{ZnO} (\ce{ZnO + 2OH- + H2O -> Zn(OH)4^2-}).
\end{solution}

\begin{solution}{exo:b1:s-block:4}
At pH 14 the water couple is \ce{2H2O + 2e- -> H2 + 2OH-}, $E = -0.83$~V;
$\log K = 2(-0.83 + 2.71)/0.059 = 64$.
\end{solution}

\begin{solution}{exo:b1:s-block:5}
In the flame, collisions raise an electron to an excited level; returning
to the \omterm{def:b1:quantum-numbers:energy-level}{ground state}, the atom emits a photon of energy equal to the gap,
specific to the element. $E = hc/\lambda = 6.626 \times 10^{-34} \times
3.00 \times 10^8/(589 \times 10^{-9}) = \qty{3.375e-19}{J} = \qty{2.11}{eV}$.
% ledger: const:h, const:c, const:eV
\end{solution}

\begin{solution}{exo:b1:s-block:6}
$10^6/106.0 = \qty{9.43e3}{mol}$ of soda ash: \qty{9.43e3}{mol} of
limestone, \qty{0.944}{t}, and $2 \times 9.43 \times 10^3/0.90 =
\qty{2.10e4}{mol}$ of salt, \qty{1.23}{t}.
\end{solution}

\begin{solution}{exo:b1:s-block:7}
$[\ce{Mg^2+}] = 1.3/24.3 = \qty{0.053}{mol/L}$; $\mathrm{pH} = 14.00 -
\frac12(11.25 + \log 0.053) = 9.0$. \qty{53.5}{mol} of magnesium in
\qty{1.0}{m^3} need \qty{53.5}{mol} of \ce{Ca(OH)2}: \qty{4.0}{kg}.
\end{solution}

\begin{solution}{exo:b1:s-block:8}
The equilibrium constant (thermodynamics) is enormous for both; the rate
(kinetics) is not decided by $E^\circ$. Lithium melts at a higher
temperature, stays solid and is covered by a less soluble product layer,
and loses its electron less easily than potassium (higher ionisation
energy): it reacts steadily, potassium violently.
\end{solution}

\begin{solution}{exo:b1:s-block:9}
$10^6/100.1 = \qty{9.99e3}{mol}$: CaO $9.99 \times 10^3 \times 56.1 =
\qty{0.560}{t}$; \ce{CO2} \qty{9.99e3}{mol}, $V = nRT/p = 9.99 \times 10^3
\times 8.314 \times 298.15/10^5 = \qty{248}{m^3}$.
\end{solution}

\begin{solution}{exo:b1:s-block:10}
$2.9 \times 10^8/6.9 = \qty{4.2e7}{kmol}$ of Li (mass in kg), $\times 73.8/2 =
\qty{1.55e9}{kg}$: about 1.6 million tonnes of lithium carbonate.
$2.9 \times 10^8/8 = \num{3.6e7}$ batteries.
\end{solution}

\begin{solution}{exo:b1:s-block:11}
1.31, 0.84, 0.71 per cent at 20, 80, \qty{100}{\celsius}: the \omterm{def:b1:precipitation:solubility}{solubility}
decreases. Hot, less carbonate stays in solution: more is recovered.
\end{solution}

\begin{solution}{exo:b1:s-block:12}
\ce{NaHCO3} is the least soluble of the salts that the ions can form, and
the solution is saturated in it first. Ammonia makes the solution basic
enough to turn \ce{CO2} into \ce{HCO3-} while being recoverable (as
\ce{NH4Cl}, regenerated by lime); sodium hydroxide would be consumed and
cost more than the product.
\end{solution}

\begin{solution}{pb:b1:s-block:1}
\textbf{1.} Li \qty{6.0}{g/L}, Mg \qty{2.4}{g/L}.
\textbf{2.} Four cubic metres of brine give one: three cubic metres of
water evaporate.
\textbf{3.} Sunlight and the dry climate supply the energy for free.
\textbf{4.} Sodium chloride, by far the most abundant salt, saturates
first.
\textbf{5.} Li $6000/6.9 = \qty{870}{mol}$; Mg $2400/24.3 = \qty{98.8}{mol}$.
\textbf{6.} Magnesium would precipitate with the carbonate and
contaminate the product.
\textbf{7.} \ce{Mg^2+ + Ca(OH)2 -> Mg(OH)2 + Ca^2+}.
\textbf{8.} $98.8 \times 74.1 = \qty{7.32}{kg}$.
\textbf{9.} $[\ce{Mg^2+}] = \num{9.88e-5}$ mol/L left: $\mathrm{pH} = 14.00 -
\frac12(11.25 - 4.01) = 10.38$.
\textbf{10.} Lithium hydroxide is soluble: no solid forms at this pH.
\textbf{11.} \ce{Ca^2+ + CO3^2- -> CaCO3}.
\textbf{12.} $98.8 \times 106.0 = \qty{10.5}{kg}$.
\textbf{13.} \ce{2Li+ + CO3^2- -> Li2CO3}.
\textbf{14.} $435 \times 106.0 = \qty{46.1}{kg}$.
\textbf{15.} $435 \times 73.8 = \qty{32.1}{kg}$.
\textbf{16.} Lithium carbonate is less soluble hot.
\textbf{17.} $0.0084 \times 1000 = \qty{8.4}{kg}$.
\textbf{18.} With small volumes of hot water, in which it is least
soluble.
\textbf{19.} $32.1 - 8.4 = \qty{23.7}{kg}$.
\textbf{20.} $23.7/32.1 = \qty{74}{\percent}$.
\textbf{21.} It still holds lithium: it is returned to the ponds.
\textbf{22.} $2.9 \times 10^8/6.0 = \num{4.8e7}$ cubic metres.
\textbf{23.} \textbf{\qty{24}{kg}} (23.7) of lithium carbonate per cubic
metre.
\end{solution}
